\section*{题目描述}

有一台神秘的多项式机器。

机器中保存着一个恰好为 $n$ 次的多项式
\[
f(x)=a_0+a_1x+a_2x^2+\cdots+a_nx^n,
\]
其中所有系数均为正整数，即
\[
a_i\in\mathbb{Z}_{>0}.
\]

你不知道这些系数的具体值。

幸运的是，这台机器允许你询问多项式在 $x=0$ 处的某阶导数。
对于一次询问，你可以给出一个整数 $k$，机器会返回
\[
f^{(k)}(0).
\]
其中约定
\[
f^{(0)}(x)=f(x).
\]

你的任务是在不超过 $n+1$ 次询问的情况下，确定多项式的所有系数
\[
a_0,a_1,\ldots,a_n.
\]

\section*{交互格式}

这是一道交互题。

交互开始时，交互器会首先输出一个整数 $n$，表示未知多项式的次数。
数据保证
\[
1\le n\le 9.
\]

\subsection*{询问}

若你想询问 $f^{(k)}(0)$，请输出：
\begin{verbatim}
? k
\end{verbatim}

其中必须满足
\[
0\le k\le n.
\]

随后，交互器会返回一个整数
\[
y=f^{(k)}(0).
\]

整个交互过程中，你最多可以进行 $n+1$ 次询问。
所有交互器返回的整数均不超过 $10^{16}$。

\subsection*{回答}

当你已经确定所有系数后，请输出：
\begin{verbatim}
! a0 a1 ... an
\end{verbatim}

其中 $a_i$ 表示 $x^i$ 的系数。
输出最终答案后，你的程序应立即结束。

\section*{交互注意事项}

每次输出询问后，都必须立即刷新标准输出缓冲区。
例如在 C++ 中可以使用：
\begin{verbatim}
cout << "? " << k << endl;
\end{verbatim}

或：
\begin{verbatim}
cout << "? " << k << '\n' << flush;
\end{verbatim}

不要尝试读取或写入任何未在交互协议中规定的信息。

如果询问次数超过限制、询问格式错误或最终答案错误，将得到 Wrong Answer。

\section*{示例}

假设机器中的多项式为
\[
f(x)=2+3x+4x^2.
\]

一种可能的交互过程如下。第一行和每次询问之后的整数由交互器发送，
以 \texttt{?} 和 \texttt{!} 开头的行由选手程序发送。
\begin{verbatim}
2
? 0
2
? 1
3
? 2
8
! 2 3 4
\end{verbatim}

这里：
\[
f(0)=2,
\]
\[
f'(0)=3,
\]
以及
\[
f''(0)=8.
\]

注意，示例仅用于说明交互格式。实际交互内容取决于隐藏的多项式。
